\section{Results}
\label{results}

% ── 1. LLM vs. Optuna head-to-head ────────────────────────────────────
%   TABLE — one row per dataset × optimizer:
%   Dataset | Optimizer | Best composite | Repro | Motif | Recall | n_sites
%   State which optimizer won per dataset and overall. Note the budget
%   asymmetry (LLM 10–12 iters, Optuna 12–15) as a caveat.
%   Multi-panel figure: iteration vs. composite for both optimizers,
%   default-parameter baseline marked. Highlights convergence speed
%   and whether LLM reasoning led to stepwise improvements.

% ── 2. Biological quality of the optimized sites ──────────────────────
%   2a. Reproducibility — bar chart per dataset: observed vs. chance-
%       expected overlap and the corrected score. Highlight RBFOX2_K562
%       as the hardest case (baseline ≈ 0.10).
%   2b. Motif enrichment — table: best motif, hit rate, background rate,
%       enrichment, canonical match (Yes/No), PWM source. Note any
%       dataset where the canonical motif is not top-ranked.
%   2c. Benchmark recall — per-dataset recovery of ENCODE reference
%       regions (chr21). Note recall vs. n_sites trade-off if visible.

% ── 3. Dataset-specific findings ──────────────────────────────────────
%   • Which parameters converged to similar values across datasets
%     (suggesting universal optima) vs. diverged (dataset-dependent)?
%   • RBFOX2_K562 vs. RBFOX2_HepG2: do optimal parameters differ for
%     the same RBP in different cell lines? (validates tuning premise)
%   • Any dataset where the LLM's reasoning trace reveals a biologically
%     interpretable decision (e.g. "widening bandwidth improved repro").

% ── 4. Collapse guard effectiveness ───────────────────────────────────
%   n_binding_sites distribution across all trials. Demonstrate healthy
%   yields under the guard; if pre-guard data exists, before/after
%   comparison for Optuna.
